\cleardoublepage

\section{绪论}

\subsection{研究背景}

\subsubsection{大语言模型与智能体的兴起}

近年来，以 \textsc{GPT}~\cite{brown2020gpt3,openai2023gpt4} 系列为代表的大语言模型（Large Language Model，LLM）在自然语言理解、逻辑推理、代码生成与多步任务规划等方面展现出了显著的涌现能力，
为人工智能从“感知智能”向“决策智能”的跨越奠定了基础。
在此基础上，智能体（Agent）成为大模型走向落地应用的关键形态~\cite{yao2023react,qin2023toollearning}：
通过赋予大模型工具调用、自主规划与多轮迭代能力，智能体能够突破自身知识边界与计算能力局限，
对接外部 \textsc{API}、专业软件与领域知识库，把自然语言意图转换成对真实系统的有效操作。
工业界提出的函数调用（Function Calling）协议~\cite{openai2023gpt4}与 \textsc{LangChain}~\cite{langchain2023}、\textsc{LlamaIndex}~\cite{llamaindex2023} 等开源框架进一步把工具封装、上下文管理与策略控制标准化，
使得构建一个可用的智能体原型仅需数十行胶水代码，整体生态正在快速成熟。

\subsubsection{科学研究场景的工具调用需求}

与此同时，AI for Science~\cite{jumper2021alphafold,butler2018ml4mat} 被广泛视为加速科学发现、提升科研效率的重要技术路径。
科学研究场景中存在大量重复性、流程化的工作，包括但不限于矩阵运算、数值积分、曲线拟合、假设检验、实验数据可视化、领域专用序列分析等，
这类工作普遍依赖 \textsc{NumPy}~\cite{harris2020numpy}、\textsc{SciPy}、\textsc{Pandas}~\cite{mckinney2010pandas} 与 \textsc{Matplotlib} 等开源科学计算库以及各类专业软件。
对于不熟悉这些工具的科研工作者而言，往往需要花费大量时间在工具学习、数据格式转换与脚本拼接上，
而真正属于“科学问题本身”的部分反而被工程性细节稀释。
因此，能够自动化调用各类专业工具、把模糊的科研需求逐步落地为可复现执行序列的科学智能体，正在成为 AI for Science 领域受到广泛关注的方向。

\subsubsection{当前智能体在科学场景下的瓶颈}

主流智能体框架在通用场景中已被广泛验证，但当我们将它们直接应用于科学研究这一更强调精度、严谨性与可复现性的垂直场景时，仍可观察到两类突出的瓶颈。
一个是工具选择困难：科学计算任务通常涉及功能高度相近、参数体系互有交集的工具，例如曲线拟合与线性回归、数值积分与微分方程求解等；
当工具库规模上升至数十乃至上百时，把全部工具描述塞入提示上下文不仅造成 token 开销随工具数量线性增长，更会显著增加 LLM 在功能相近工具间的选择错误率~\cite{qin2023toollearning}，
而这一问题在朴素的 ReAct 主循环~\cite{yao2023react}下并未得到结构化的处理。
另一个是任务拆解薄弱：科研任务的常见形态是“矩阵运算 $\rightarrow$ 求方程 $\rightarrow$ 数值积分 $\rightarrow$ 统计分析 $\rightarrow$ 绘图”一类的长依赖链，
其中每一步的输入都依赖前序步骤的具体输出字段。
朴素的 ReAct 与单纯的自然语言 Plan-and-Solve~\cite{wang2023plansolve} 仅依靠 LLM 自身的自由生成来组织步骤，缺乏一种纯代码可校验的结构化规划支撑，
极易出现参数遗漏、依赖错绑、占位符未解析等问题，难以满足科研工作者对“每一步都可追溯、可复现、可解释”的现实诉求。

此外，科学场景对中间数值的精度、单位、符号约定等都有严格要求。
通用 LLM 在面对 $R^2$ 拟合优度为负、统计量量级失配、计算结果中出现 \texttt{NaN} 等异常情形时，
往往会用“看起来合理”的字面回答掩盖底层错误；
若不在系统侧引入针对科研场景的合理性检查与精度规范，工具调用结果的可靠性就难以保障。

\subsection{研究目标}

针对上述瓶颈，本文以\textbf{面向科学场景的智能体工具调用优化}为核心研究问题，
试图在不依赖额外人工监督数据、不依赖特定基座模型微调的前提下，
仅通过算法与工程层面的设计提升智能体在科学任务上的工具选择精度、拆解严谨性、调用容错性与轨迹可解释性。
具体而言，本文将研究目标聚焦为以下三层。

算法层面，提出两个可独立消融、形式化清晰的核心算法模块——
检索增强工具选择算法（Retrieval-Augmented Tool Selection，RATS）与
结构化 DAG 任务拆解算法（DAG Planner），
分别对应“选什么工具”与“怎么组织工具调用”这两个朴素 ReAct 范式中尚未被结构化处理的问题。

工程层面，在算法模块周边构建一套可对外讲述的科研基础设施，
包括基于统一抽象基类的标准化科学计算工具库、面向科学场景的提示工程规则集，
以及覆盖三维判分（数值正确性、文件正确性、工具覆盖性）的端到端评估框架。

实验层面，构建一个能区分不同方法贡献度的科学任务测试集与多基线对比框架，
通过同口径的全量评测从成功率、迭代次数、工具调用次数与端到端耗时四个维度出发，
对方法的有效性、效率与失败模式进行透明、可复现的实证分析。

\subsection{主要工作与贡献}

围绕上述目标，本文的主要工作与贡献如下。

\paragraph{检索增强工具选择算法（RATS）}
提出一种结合嵌入语义检索与规则化重排的轻量工具选择算法。
RATS 通过离线预计算工具描述向量、在线对任务向量做余弦相似度检索，并融合嵌入相似度、参数重叠度、历史成功率三类特征进行规则化重排，
在保留最低候选数的保底机制下，将传给 LLM 的工具上下文从全量 $N$ 压缩到精简的 $K$。
该算法不需要任何额外训练，可以与任意支持函数调用接口的 LLM 自由组合，
本质上把检索增强生成（Retrieval-Augmented Generation）的范式迁移到了“工具”这一新的检索对象。

\paragraph{结构化 DAG 任务拆解算法（DAGPlanner）}
设计一种基于有向无环图（DAG）模式的结构化规划器，
通过显式定义支持依赖占位符 \texttt{\$\{sX.field\}} 的计划模式（Plan Schema），让 LLM 一次性产出可被纯代码校验的执行计划。
配套的 \texttt{PlanValidator} 在不调用 LLM 的前提下覆盖五类失败模式——模式不符、工具名不在候选集、依赖图有环、变量未声明、占位符越界，
并以错误驱动的方式触发携带诊断信息的有限次重规划；
当重规划在硬上限内仍未收敛时，系统将平滑降级到 ReAct 主循环作为安全网，从而在结构化能力与鲁棒性之间取得均衡。

\paragraph{科学场景提示工程规则集}
提炼出 4 条针对科学场景的提示工程规则（O1 结果合理性检查、O2 数值计算规范、O3 开场规划模板、O4 Few-shot 示例注入），
作为 RATS 与 DAGPlanner 在运行时不可或缺的辅助约束。
通过同口径对比实验证明，这 4 条规则与算法模块呈现明显的协同效应，
共同构成本文方法相较于业界基线的可解释优势来源。

\paragraph{评估框架与基准测试集}
构建覆盖数值计算、统计分析、可视化、组合任务、错误恢复 5 大类别共 54 道科学任务的测试集，
并实现一个支持三维判分、按配置/按用例任意组合复现的评测框架。
在该框架下完成了包括 CoT、Plan-and-Solve、Reflexion 在内 9 个配置 $\times$ 54 题共 486 次运行的全量评测，
为后续从主实验、按类别细分、消融与效率四个维度展开论证提供了同口径的实证基础。

\subsection{论文组织结构}

本文余下章节按如下方式组织。
第二章对智能体框架、任务规划、工具检索与选择以及检索增强生成四个相关研究方向进行综述，并指出本文方法与其差异点。
第三章是论文的核心算法章，给出问题形式化定义后，分别详述 RATS 与 DAGPlanner 的设计动机、形式化流程、复杂度与收敛性分析，并说明二者如何集成为本文的完整方案 \texttt{Ours\_full}。
第四章介绍对应的系统实现，包括标准化科学计算工具库、Agent 框架、提示工程体系与评估框架。
第五章围绕 9 配置 $\times$ 54 题的全量评测，从主实验、类别细分、消融与效率四个角度对方法进行实证分析。
第六章聚焦关键案例与失败模式，对若干反直觉现象（如 CoT 反而降分）进行深度讨论。
第七章总结全文工作，并就方法局限性与未来工作做出诚实的表述。
